\documentclass[aspectratio=169,11pt]{beamer}
\usetheme{default}
\usecolortheme{seahorse}
\setbeamertemplate{navigation symbols}{}
\title{Q13: How Models Represent and Store Knowledge}
\author{Paradise (nuc13)}
\date{October 2026}
\begin{document}
\frame{\titlepage}

\begin{frame}{The question}
\begin{itemize}
\item ``How do models represent and store knowledge, and use it to answer?''
\item No database, no schema --- yet facts come out. Where do they live?
\item Answer preview: \textbf{mid-layer FFN modules as key--value memories; features as directions in superposition; retrieval routed by attention}
\end{itemize}
\end{frame}

\begin{frame}{Storage site: FFN layers are key--value memories}
\begin{itemize}
\item FFN = $\approx$ two-thirds of Transformer parameters
\item $\mathrm{FFN}(x) = f(xK^\top)\,V$ \quad $\Rightarrow$ structurally a key--value memory (Geva et al.\ 2021)
\item Keys: rows of $K$ match textual patterns; values: columns of $V$ induce output distributions
\item Lower layers: shallow patterns; upper layers: semantic themes
\item Output = weighted \emph{composition} of memories, refined via residual stream
\end{itemize}
\end{frame}

\begin{frame}{Retrieval path: causal tracing (ROME)}
\begin{itemize}
\item Intervention: corrupt the input; restore states one at a time; see which repairs the fact
\item Finding (GPT-2 XL): mid-layer \textbf{MLP modules} at the \textbf{last subject token} mediate recall
\item Late-layer \textbf{attention} then moves the information to the output position
\item Storage (MLP) and routing (attention) are \emph{separable} mechanisms
\end{itemize}
\end{frame}

\begin{frame}{Format: linear representations}
\begin{itemize}
\item Linear representation hypothesis: features = \textbf{directions}; activations = $\sum_i a_i f_i$
\item Why probing works: linear readout suffices --- if information is used or at least accessible
\item Caveat: probing accuracy $\neq$ causal use; complement with interventions
\item But then: a $d$-dimensional layer could only hold $d$ features\dots
\end{itemize}
\end{frame}

\begin{frame}{Format: superposition}
\begin{itemize}
\item Networks store \textbf{more features than dimensions} (Elhage et al.\ 2022)
\item Trick: non-orthogonal, nearly-collapsed directions; high dimensions make random vectors almost orthogonal
\item Cost: \textbf{interference} --- co-active features leak into each other's readout
\item ReLU acts as the noise filter; geometric structures (triangles, simplices) emerge
\item \textbf{Polysemanticity} is the visible symptom --- one neuron, many concepts
\end{itemize}
\end{frame}

\begin{frame}{Demo: superposition in NumPy}
\begin{itemize}
\item $N$ features into $d < N$ dimensions via random near-orthogonal vectors
\item Sparse features: recoverable readout despite interference
\item Dense features: interference dominates --- capacity exceeded
\item \texttt{code/superposition\_demo.py} --- pass/fail assertions, no dependencies beyond NumPy
\end{itemize}
\end{frame}

\begin{frame}{Editing: locate, then patch}
\begin{itemize}
\item \textbf{ROME}: rank-one update of one mid-MLP matrix; $k_* \mapsto v_*$
\item Generalizes across paraphrases, stays specific to the subject
\item \textbf{MEMIT}: spreads batches of edits across mediating layers (thousands at once)
\item Side effects: bleedover to neighbors, fluency degradation, brittleness
\item Open problem: surgical edits in a superposed store, with guarantees
\end{itemize}
\end{frame}

\begin{frame}{In weights vs.\ in context}
\begin{itemize}
\item Weights: vast, static, hard to update (edit/retrain)
\item Context: tiny, instant, attention is the read mechanism
\item Failure modes differ: stale/absent vs.\ retrieval error / overflow
\item Forward-link to Q12: what happens when the two channels \textbf{conflict}
\end{itemize}
\end{frame}

\begin{frame}{Reading}
\begin{itemize}
\item Geva et al.\ 2021 --- Transformer Feed-Forward Layers Are Key-Value Memories (EMNLP)
\item Meng et al.\ 2022 --- Locating and Editing Factual Associations in GPT (NeurIPS)
\item Elhage et al.\ 2022 --- Toy Models of Superposition (Transformer Circuits)
\item Meng et al.\ 2023 --- Mass-Editing Memory in a Transformer (ICLR)
\end{itemize}
\end{frame}

\end{document}
